\section{Practice Exercise}
\begin{practiceexercise}
\textbf{Multiple Choice Questions}
\begin{enumerate}
  \item Which of the following statements represents an authentic expression of a need rather than a strategy or pseudo-need \citep{gordon1970}?
    \begin{itemize}
      \item[A)] ``I need you to take out the trash right now.'' \citep{cohen2001}
      \item[B)] ``I need you to be more responsible.'' \citep{kohn2005}
      \item[C)] ``I need to feel supported and have a sense of fairness in our household chores.''
      \item[D)] ``I need you to stop playing video games all evening.''
    \end{itemize}
  \vspace{\fill}
  \item Your teenager comes home angry, slamming the door, and says, ``You always treat me like a baby!'' Which underlying need is most likely driving this outburst \citep{brown2012}?
    \begin{itemize}
      \item[A)] A need for you to stop talking to them. \citep{dweck2006}
      \item[B)] A need for autonomy, trust, and to be respected as a young adult.
      \item[C)] A need for more rules to guide their behavior.
      \item[D)] A need for an apology for ruining their day.
    \end{itemize}
  \vspace{\fill}
  \item When translating a partner's judgment (e.g., ``You're so selfish!'') into a need, what is the best empathetic response \citep{bowlby1969}?
    \begin{itemize}
      \item[A)] ``I am not selfish; I do plenty for this family.''
      \item[B)] ``Are you feeling upset because you need to feel considered and valued?''
      \item[C)] ``Well, you're selfish too sometimes.''
      \item[D)] ``You need to stop criticizing me.''
    \end{itemize}
\end{enumerate}
\vspace{\fill}
\newpage

\textbf{Explanations for Questions 1--3}

\textbf{Question 1: C.} This statement clearly articulates universal human desires (support, fairness) without attacking the other person's character or demanding a specific, controlling action \citep{dunn2000}. Options A and D are strategies disguised as needs \citep{kabatzinn1990}, and Option B is a judgment disguised as a need \citep{rogers1961}.

\textbf{Question 2: B.} Anger and accusations of being treated ``like a baby'' typically stem from a teenager's thwarted desire for independence and autonomy \citep{thomas1977}. Recognizing this allows the parent to bypass the defensive tone and address the core issue of trust and growing independence \citep{lazarus1991}. A failure to address this can limit proper emotional coaching \citep{gottman1997}.

\textbf{Question 3: B.} This response perfectly demonstrates empathetic translation \citep{siegel2011}. It bypasses the judgmental attack (``selfish'') and checks in with the partner's underlying feeling (upset) and core need (to be considered/valued) \citep{rosenberg2003}, inviting connection rather than starting an argument \citep{greene1998}.
\vspace{\fill}
\newpage
\end{practiceexercise}
